% PRIVATE. Research document copied 2026-10-08 into spine-program/09_wrong_level_clinical_qi from
% spinesurg-ct-nnunet (private, branch L6_Class, commit 097933a) paper/qi_universal_protocol/main.tex
% (last commit d4b895e, 2026-09-22). It describes the sweep and the persistent digital fiducial, which are
% invention aspects: it stays private until Greg confirms the provisional is filed (HANDOFF MANUSCRIPTS.md
% section 5). Known defects, left as written in this copy: the affiliation line reads Neurosurgery and must
% read General Surgery (Wayne State University School of Medicine / Detroit Medical Center); 'the encoder
% takes 17 ms per frame' has no result file; 'about fifteen minutes for fluoroscopic localization' is an
% unverified citation. See TODO.md step 4. Edits made in this copy: the corresponding e-mail is withheld;
% \bibliography points at refs.bib beside this file (a copy of paper/rsnai/refs.bib at 180fb0f).
% Wrong-level spine surgery after the Universal Protocol: what the evidence supports, and a quality-
% improvement protocol for a mechanism that names the level rather than re-counting it.
% Target: The Joint Commission Journal on Quality and Patient Safety (analysis / QI protocol; SQUIRE 2.0
% for the protocol section). Every figure traces to docs/WRONG_LEVEL_COST_EVIDENCE.md and
% docs/WRONG_SITE_BURDEN_EVIDENCE.md, where each carries its adversarial-verification vote; derived
% arithmetic is labelled as derived in the text.
\documentclass[11pt,letterpaper]{article}
\usepackage[margin=1in]{geometry}
\usepackage[T1]{fontenc}
\usepackage{lmodern}
\usepackage{helvet}
\renewcommand{\familydefault}{\sfdefault}
\usepackage{setspace}
\usepackage{booktabs}
\usepackage{natbib}
\usepackage{textcomp}
\setstretch{1.15}

\begin{document}

\noindent{\Large\bfseries The Universal Protocol with a level in it: site marking and time-out for spine
surgery assisted by a fluoroscopic level map---evidence, and a quality-improvement protocol}

\bigskip
\noindent Gregory Schwing, MD, PhD\\
Department of Neurosurgery, Wayne State University School of Medicine / Detroit Medical Center\\
\texttt{[e-mail withheld]}

\bigskip
\noindent\textbf{Word count:} about 2,400 (main text). \textbf{Tables:} 2.

\section*{Abstract}
\noindent\textbf{Background.} The Universal Protocol has been mandatory in accredited United States
hospitals since 2004. Wrong-level spine surgery persists under it, and the reason is structural: a
time-out verifies a site that has been \emph{marked}, but a vertebral level is not marked, it is counted,
on intraoperative images that show a few bodies at a time.

\noindent\textbf{Evidence.} Two adversarially verified literature syntheses (every claim put to three
independent refutation attempts) give the following. Half of responding neurosurgeons report at least one
wrong-level operation in their career. Per-procedure incidence is 0.032\,\% by self-report and 0.12\,\% in
a prospective 34,544-case registry, with the thoracic spine highest at 0.42\,\%. Chart review found the
Universal Protocol would have prevented 62\,\% of insured wrong-site cases under ideal conditions, and
wrong-site events continued at a steady rate through the protocol era. In a contemporary closed-claims
series the mean wrong-site indemnity was \$136,453 and spine surgery the commonest procedure; 17\,\% of
self-reported wrong-level events led to legal action and 81\,\% of litigated cases ended in payment. No
published study measures the intraoperative localisation \emph{process}---exposures, minutes, dose, needle
counts, disagreements---as a quality metric, and no per-event hospital cost exists.

\noindent\textbf{Proposal.} A single-centre, SQUIRE-structured quality-improvement study of the Universal
Protocol performed \emph{with} a level in it. The protocol's three elements map onto a fluoroscopic level
map built in theatre: verification enters the target level; site marking is a short sweep of overlapping
exposures from the sacrum, linked by software into one map, that places a digital fiducial on the target
body in the live image; the time-out confirms the fiducial's level against the surgeon's own count before
incision, and a disagreement is resolved by rule. On held-out CT renders the walk from the sacrum places
the fiducial on the right body in 94.5\,\% of 950 walks (96.3\,\% of those reached) from two exposures on
average, and 92.2\,\% of 7,405 walks on 522 clinical CTs from a second institution; published C-arm
repositioning times put the whole manoeuvre at one to three minutes against fifteen for fluoroscopic
counting. After a four-month process baseline, the assisted protocol runs for eight months; the primary
outcome is the disagreement rate between the surgeon's count and the fiducial, adjudicated against
postoperative imaging.

\noindent\textbf{Conclusion.} The Universal Protocol is necessary and not sufficient for wrong-level
spine surgery, because it verifies a mark and the spine has none. Giving it a mark that was named from
the patient's own images, and a time-out that checks that name against the count, is a change to the
protocol rather than a replacement for it, and it can be measured at one centre in a year.

\section*{The problem the Universal Protocol was not built for}
The Universal Protocol for Preventing Wrong Site, Wrong Procedure, and Wrong Person Surgery requires
preprocedure verification, marking of the procedure site, and a time-out; for spinal procedures it
allows that intraoperative imaging may be needed to locate and mark the exact vertebral level
\citep{jcnpsg2025}. Each element
assumes the site can be identified on the patient before incision. For a laterality error that is true:
the limb is visible and the mark sits on it. For a vertebral level it is not. The level is identified
intraoperatively on fluoroscopy, by placing needles every few levels and counting cephalad from the
sacrum on lateral views, or by counting levels on anteroposterior views \citep{strong2021,mayer2008}. The time-out can confirm that everyone agrees on the target level; it
cannot confirm that the count reached it.

Two independent surveys found that about half of responding neurosurgeons had performed at least one
wrong-level spine operation \citep{mody2008,groff2013}. Self-reported incidence was one in 3,110
procedures \citep{mody2008}; a prospectively collected multicentre registry of 34,544 degenerative spine
operations, in which intraoperative disorientation caused 81\,\% of the lumbar events, recorded 0.12\,\% overall,
0.42\,\% thoracic, 0.13\,\% lumbar and 0.05\,\% cervical \citep{furuya2026}. A systematic review found
twelve incidence studies, none thoracic-specific, and reported that in the one study assessing it,
site-verification protocols failed to prevent about one-third of cases \citep{longo2012}. The
frequently quoted 0.09--4.5 per 10,000 range is wrong-site surgery across all specialties, not a spine
rate \citep{devine2010}.

\section*{What the protocol era shows}
Kwaan and colleagues reviewed 13 insured nonspine wrong-site cases from 1985--2004 and judged
retrospectively that the Universal Protocol, under optimal conditions, might have prevented eight
(62\,\%) \citep{kwaan2006}. Stahel and colleagues
analysed 27,370 physician self-reported occurrences in a Colorado insurer database from 2002 to mid-2008
and found 107 wrong-site and 25 wrong-patient procedures, persisting after the Universal Protocol was
introduced, with significant harm in 35.5\,\% of wrong-site cases and one death \citep{stahel2010}. The Joint Commission received 127 reports of wrong surgery in
2024, nine wrong-level, seven of those spinal, from a voluntary system that, it cautions, cannot show
how often such events actually occur \citep{jc2024}. NHS England provisionally recorded 185 wrong site surgery Never Events in
2024/25, a category that also counts wrong-site blocks and wrong skin lesions \citep{nhse2025}. None of these series can say how many wrong-level spine operations occur; all of them
say the number has not gone to zero.

The costs that are measurable are legal. Seventeen percent of self-reported wrong-level events led to
legal action or settlement \citep{mody2008}; of 68 litigated wrong-disc-level cases, 55 ended in a
payment and surgeons prevailed in 42\,\% of jury trials \citep{goodkin2004}; the mean wrong-site
indemnity in a 2013--2020 closed-claims series was \$136,453 (median \$36,250), with spine and disc
surgery the commonest procedure and 45.6\,\% of claims involving additional surgery \citep{tan2023};
malpractice payments for all surgical never events summed to \$1.3 billion over 21 years
\citep{mehtsun2013}; 978 spinal-surgery negligence claims in England over five years were valued at
\pounds535.5 million \citep{machin2018}. Applying the two per-procedure rates to 1.3--1.6 million spine
operations a year gives, as derived arithmetic, roughly 400--1,900 wrong-level operations and \$8--37
million in indemnity a year in the United States; the hospital cost of the reoperations, added stay and
readmissions those events cause has no published per-event figure, and there is no defensible worldwide
estimate.

\begin{table}[htbp]\centering
\caption{The verified figures and where each comes from. Votes are refutation attempts survived (of 3).}
\small
\begin{tabular}{p{7.2cm}p{3.2cm}p{4.8cm}}
\toprule
Figure & Source & Verification \\
\midrule
Half of neurosurgeons report a wrong-level operation & Mody 2008; Groff 2013 & 3/3 twice \\
Incidence 1 in 3,110 (self-report); 0.12\,\% (registry) & Mody 2008; Furuya 2026 & 3/3 twice \\
Thoracic 0.42\,\%, lumbar 0.13\,\%, cervical 0.05\,\% & Furuya 2026 & 3/3 \\
Universal Protocol would have prevented 62\,\% & Kwaan 2006 & 2/3 \\
Steady rate through the protocol era & Stahel 2010 & 3/3 \\
17\,\% of events litigated; 81\,\% of litigated paid & Mody 2008; Goodkin 2004 & 3/3, 2/3 \\
Mean wrong-site indemnity \$136,453; spine commonest & Tan 2023 & 3/3 \\
\$1.3 billion never-event payments, 21 years & Mehtsun 2013 & 3/3 \\
179 fiducial placements, 51.9 min, 183 mGy & AJNR 2023 & 3/3 \\
Operating-room time \$34--37 per minute & Childers 2018 & 3/3 \\
\bottomrule
\end{tabular}
\end{table}

\section*{Why counting fails, and what it costs to avoid counting}
The traditional intraoperative technique is needles every few levels, counted from the sacrum on
lateral fluoroscopy; surgeons disagree even on which structure to count, the facet joint with its
pedicle for 67\,\% and the spinous process for 49\,\% in the thoracic spine \citep{strong2021,mayer2008}. The images supply the failure
modes. A field of a few bodies means a count to a mid-thoracic level passes through several
exposures, each of which can be joined to the previous one body off. A lumbosacral or thoracolumbar
transitional vertebra shifts the count for the whole column, and such variants are common enough that
any series of a few hundred spines contains them. Ribs are hard to number where a twelfth rib is absent
or hypoplastic. Habitus, osteopenia and the scapular shadow degrade the mid-thoracic image most. None of
these is a lapse of attention that a time-out addresses.

The cost of avoiding the count is visible at one centre, Mayo Clinic in Rochester, which placed
preoperative percutaneous gold fiducial markers in 179 patients over 38 months, 176 of them thoracic,
at a mean of 51.9 minutes and 183\,mGy each, to let the surgeon find the level \citep{siminski2023}. At a
hospital operating-room cost of \$34--37 per minute before physician fees \citep{childers2018}, each
placement represents about \$1,750--1,950 of staffed procedural time and the series about 155 hours; the
figure is a proxy, since the marker is placed in a fluoroscopy or CT suite, and it omits the second
anaesthetic and the second day. That a centre with every resource chooses this over counting is the
clearest statement in the literature that the thoracic level is, in practice, un-nameable in theatre.

\section*{The measurement gap}
The literature measures outcomes---events, claims, payments---and does not measure the process that
produces them. No published study records, per case, how many fluoroscopic exposures were spent on
localisation, how many minutes, what dose, how many needles, whether the surgeon and the radiographer
disagreed, or how a disagreement was resolved. Without that baseline no intervention can be evaluated,
and the one intervention the field has, the time-out, was evaluated retrospectively against claims. A
quality-improvement study that measures the process is the missing first step, and it is cheap.

\section*{What the engine does, and what it has been measured at}
The level map is built from what the radiographer already does. The first exposure places the sacrum at
the bottom of the field; each move is one field minus two vertebral bodies, and the last move is a full
move or none. Software links each new frame to the last by image correlation restricted to the step the
protocol asked for, so the periodic column cannot lock one body off, and names every body in the map by
counting from the sacrum and by a shape reader, with the two readings reconciled. When the target level
enters a frame the map marks it: a digital fiducial drawn on the live image, on the body the map says is
the target. That is the site mark the Universal Protocol asks for and the spine has never had.

Measured on renders of CT the software was never trained on, walking from the sacrum to every labelled
level in turn, the fiducial landed on the right body in 94.5\,\% of 950 walks on 120 spines (96.3\,\% of
the walks that reached the target) from 2.0 exposures on average, and in 92.2\,\% of 7,405 walks on 522
clinical CTs from another institution, 96.2\,\% in the lumbar and 91.6\,\% in the thoracic spine, from
3.0 exposures. The errors are one body off and are concentrated in the known ambiguities: a spine with
four lumbar bodies on the sacrum, or six, where a count from one anchor cannot tell one more lumbar body
from one fewer thoracic; the decoder has since gained the four-lumbar alphabet and the remainder are
stated as the residual. Accuracy falls with distance from the anchor, 96\,\% within six levels and about
90\,\% at ten to sixteen, which is why thoracic work is anchored from C2 or the first rib rather than
walked twenty bodies from the sacrum. On 25 real lumbar radiographs, with the detector's own boxes, the
map completed on 21; the film domain is where the residual risk is, and it is why the study below keeps
the surgeon's count as the decision.

The time is the C-arm's, not the software's: the encoder takes 17 ms per frame and the linking a fraction
of a second per walk, while published timed positioning trials put reaching a target view with a C-arm at 36
to 49 seconds on average \citep{matthews2007,desilva2018}. A lumbar walk of two exposures is therefore about a minute and a
mid-thoracic walk of three to four about two to three, against about fifteen minutes for fluoroscopic
localization in the thoracic spine and about three with a pre-placed fiducial in one small
non-randomised series \citep{mongardi2021}. The fiducial moves the work, not removes it: a separate CT
procedure of a mean 14.3 minutes 2--48 hours before surgery \citep{keil2024}, usually as an outpatient
the day before \citep{siminski2023}.

\section*{A quality-improvement protocol (SQUIRE 2.0)}
\noindent\textbf{Setting and population.} One academic neurosurgery service; all elective and urgent
thoracic and lumbar operations requiring intraoperative level localisation over twelve months, with
cervical cases as a comparison arm. Institutional review as quality improvement.

\noindent\textbf{Phase 1, baseline (months 1--4).} For every case the circulating radiographer records
on a one-page form: the anchor used (sacrum, C2, ribs, preoperative marker), the number of localisation
exposures, fluoroscopy minutes and cumulative dose from the C-arm log, the number of needles, whether a
preoperative fiducial or CT-guided marker was placed (with its procedure time and dose from the
radiology record), the surgeon's stated confidence before incision, and any disagreement between team
members and how it was settled. The level reached is adjudicated against postoperative imaging by a
radiologist blinded to the form. Transitional anatomy is read from the preoperative CT where one exists.

\noindent\textbf{Phase 2, the assisted protocol (months 5--12).} The Universal Protocol is performed as
written, with each element carried by the map. \emph{Verification:} the target level from the consent
and the plan is entered into the system before the patient is draped. \emph{Site marking:} at the start of
localisation the radiographer performs the sweep from the sacrum (or from C2 for cervical and upper
thoracic work) until the system reports the target in the field, and the digital fiducial is displayed on
the live image and on the monitor the team sees. \emph{Time-out:} the surgeon states the level by their
own count, the circulator reads the fiducial's level and the system's confidence, and the two are
compared aloud with the plan. Agreement proceeds. Disagreement, or a system refusal, invokes a written
rule: one more exposure with greater overlap, then a count from the other anchor if one is in reach, then
the surgeon's judgement, recorded as such. The surgeon's count remains the decision throughout; the map
adds a second, independent reading and a mark. Every disagreement is recorded with its resolution and the
adjudicated truth. The system's own accuracy is published separately and is not an outcome here.

\noindent\textbf{Outcomes.} Primary: the disagreement rate between the surgeon's count and the fiducial
at the time-out, and which of the two the adjudicated level supported, by region. Compliance: the share
of cases in which all three elements were performed as written, and the share in which the fiducial was
in the field at incision. Secondary: localisation exposures,
minutes and dose per case, needles per case, preoperative marker placements per month, surgeon
confidence, and the near-miss rate (a count corrected before incision). Balancing: total operative
time, and any case in which the map was wrong and the surgeon right.

\noindent\textbf{Analysis.} Phase 1 gives the process distribution per region, which is the number this
literature lacks. Phase 2 is analysed as an interrupted series on exposures and minutes, and the
disagreement rate is reported with its interval and every case listed. The study is powered to
estimate, not to detect a change in wrong-level events, which at 0.12\,\% would need tens of thousands
of cases; the near-miss rate is the tractable proxy and is reported as such.

\begin{table}[htbp]\centering
\caption{Process measures recorded per case, and the source of each.}
\small
\begin{tabular}{p{6.5cm}p{4.2cm}p{4.5cm}}
\toprule
Measure & Source & Why \\
\midrule
Anchor used & radiographer form & which count was made \\
Localisation exposures, minutes, dose & C-arm log & the cost of counting \\
Needles placed & form & the count's length \\
Preoperative marker (time, dose) & radiology record & the cost of not counting \\
Surgeon confidence before incision & form (5-point) & when the count is doubted \\
Disagreement and resolution & form & the near-miss \\
Level reached & postoperative imaging, blinded & the truth \\
Transitional anatomy & preoperative CT & the known failure mode \\
\bottomrule
\end{tabular}
\end{table}

\section*{Limitations}
The evidence synthesis is of published figures, several of them self-reported or claims-based, and the
annual event and dollar figures are derived from stated rates and volumes rather than counted. The
proposed study is single-centre and cannot measure a change in wrong-level events; it measures the
process and the near-miss. The second reader is one implementation of a level map, and its own accuracy
is a separate question with separate evidence.

\section*{Conclusion}
The Universal Protocol verifies a mark. A vertebral level has no mark; it has a count, made on images
that show a few bodies at a time, through anatomy that varies. That is why wrong-level spine surgery
has persisted through two decades of time-outs, why a centre pays an hour and a second procedure per
thoracic patient to avoid counting, and why the next quality-improvement step is not another checklist
but a measurement of how the level is named today and a second reader that names it independently.

\bibliographystyle{unsrtnat}
\bibliography{refs}
\end{document}
